%% ==================================================
%% Section: The tree circuit with measurements
%% Sources: arXiv:2307.01696 (Malz, Styliaris, Wei, Cirac), eq. (16) and
%% paragraph "Tree-RG circuit with measurements"; arXiv:2103.13367 (Piroli,
%% Styliaris, Cirac), paragraphs "State transformations with QC and LOCC" and
%% "Phases of matter", and Example 2.
%% ==================================================

\section{The tree circuit with measurements}\label{sec:ldp_tree_measurements}

The isometries of the tree circuit act on registers that are far apart. They
``can be teleported at neighboring registers with a constant overhead''
\cite{Malz2024Preparation}, so that every layer of the tree takes constant
depth. Theorem~\ref{thm:qc_long_range_layer} applies a layer of two-site gates
joining the sites $a$ and $a+2L+1$, on pairwise disjoint stretches
$\{a,\dots,a+2L+1\}$, with measurements in depth $5$ for every $L$. This
section deduces that a binary tree of two-site gates with $k$ levels is
applied with measurements in depth $5k$. Every register here is one site of
the ring. The registers of the tree circuit carry $\mathbb C^{D^2}$; registers
of several sites, and the resulting depth $O(\log\log(N/\varepsilon))$ for
matrix product states, are the subject of Section~\ref{sec:ldp_register_tree}.

Throughout, the sites of the ring are $\{0,\dots,N-1\}$ with addition modulo
$N$, and the local labels are $\mathbb Z_d=\{0,\dots,d-1\}$ with $d\ge1$.

\begin{lemma}[Product vectors with zero at a set of sites]
    \label{lem:ldp_zero_sites_product_vector}
    \lean{MPSPreparation.isZeroOn_productVector}
    \leanok
    \uses{lem:qc_site_permutations}
    A product vector whose site vectors at the sites of a set $Z$ are multiples
    of $\ket0$ has $\ket0$ at the sites of $Z$, in the sense of
    Lemma~\ref{lem:qc_site_permutations}.
\end{lemma}
\begin{proof}
    \leanok
    The product vector vanishes at a configuration $x$ as soon as one factor
    $\pi_i(x_i)$ with $i\in Z$ and $x_i\neq0$ vanishes.
\end{proof}

\begin{theorem}[Binary trees of gates in depth $5k$ with measurements]
    \label{thm:ldp_tree_measurements}
    \lean{MPSPreparation.blockInterior, MPSPreparation.treeInterior,
        MPSPreparation.treeInterior_mono,
        MPSPreparation.blockInterior_subset_treeInterior, MPSPreparation.treeGate,
        MPSPreparation.levelGates, MPSPreparation.treeGate_a_val,
        MPSPreparation.treeGate_far_val, MPSPreparation.pairwise_levelGates,
        MPSPreparation.levelOp, MPSPreparation.treeOp,
        MPSPreparation.treeRounds, MPSPreparation.sum_depth_treeRounds,
        MPSPreparation.isZeroOn_levelOp_mulVec,
        MPSPreparation.interior_subset_blockInterior,
        MPSPreparation.isRoundsImplementationOn_treeRounds,
        MPSPreparation.isPreparedWithMeasurementRoundsInDepth_treeOp}
    \leanok
    \uses{thm:qc_long_range_layer, lem:qc_round_implementation_compose,
        lem:ldp_zero_sites_product_vector}
    For $i\ge0$, level $i$ of the tree applies, for every block
    $\{p2^{i+1},\dots,(p+1)2^{i+1}-1\}$ with $(p+1)2^{i+1}\le N$, a two-site
    unitary $u_{i,p}$ to the two endpoints $p2^{i+1}$ and
    $p2^{i+1}+2^{i+1}-1$ of the block. Let $T_k$ be the product of the levels
    $0,\dots,k-1$, the level $k-1$ applied first. Let $I_B$ be the set of
    sites $x$ whose residue modulo $B$ lies in $\{1,\dots,B-2\}$ and whose
    block $\{qB,\dots,(q+1)B-1\}$, $q=\lfloor x/B\rfloor$, satisfies
    $(q+1)B\le N$, and let $J_k=I_2\cup I_4\cup\dots\cup I_{2^k}$. Then $2k$
    rounds of total depth $5k$ implement $T_k$ on the vectors with $\ket0$ at
    the sites of $J_k$. In particular, for a nonzero product vector $\ket\pi$
    whose site vectors on $J_k$ are multiples of $\ket0$, the vector
    $T_k\ket\pi$ is prepared with measurement rounds in depth $5k$.
\end{theorem}
\begin{proof}
    \leanok
    The gate of the block $p$ of level $i$ is a distant gate with $a=p2^{i+1}$
    and $L=2^i-1$. Its stretch is the block, so the stretches of a level are
    pairwise disjoint, and its interior lies in $I_{2^{i+1}}\subseteq J_{i+1}$.
    By Theorem~\ref{thm:qc_long_range_layer} two rounds of depth $5$ implement
    level $i$ on the vectors with $\ket0$ at $J_{i+1}$. The gates of level $i$
    act on sites $x$ with $x\equiv0$ or $x+1\equiv0$ modulo $2^{i+1}$, hence
    modulo $2^{j+1}$ for every $j<i$, so on no site of $J_i$; and
    $J_i\subseteq J_{i+1}$. So level $i$ maps the vectors with $\ket0$ at
    $J_{i+1}$ to vectors with $\ket0$ at $J_i$, and
    Lemma~\ref{lem:qc_round_implementation_compose} composes the levels by
    induction over $k$; $J_0$ is empty.
\end{proof}

\begin{remark}[The tree circuit of matrix product states]
    \label{rem:ldp_tree_measurements_mps}
    \uses{thm:ldp_tree_measurements, thm:ldp_loglog_depth,
        thm:ldp_loglog_depth_every_length}
    In the tree circuit (\ref{eq:ldp_tree_factorization}) a register above
    the lowest level carries $\mathbb C^{D^2}$, which is a group of sites once
    $D^2>d$. Theorem~\ref{thm:ldp_loglog_depth} takes registers of $s$
    consecutive sites, $s$ a length at which the blocked tensor is injective,
    and obtains the depth $O(\log\log(N/\varepsilon))$ of
    \cite{Malz2024Preparation} for normal tensors and the chain lengths
    divisible by $s2^{k+1}$ for some $k$ with
    $a\log(N/\varepsilon)+b\le s2^{k+1}\le2(a\log(N/\varepsilon)+b)$, and
    Theorem~\ref{thm:ldp_loglog_depth_every_length} obtains it for every chain
    length, with trees on leaves of unequal widths
    \cite{gap:mswc24_tree_measurement_scope}.
\end{remark}
